\documentclass[11pt,a4paper]{article}
\usepackage[margin=2.4cm]{geometry}
\usepackage{amsmath,amssymb}
\usepackage{booktabs}
\usepackage{longtable}
\usepackage{array}
\usepackage[T1]{fontenc}
\usepackage{lmodern}
\usepackage{microtype}
\usepackage{xcolor}
\usepackage[colorlinks=true,linkcolor=blue!50!black,urlcolor=blue!50!black]{hyperref}

\newcommand{\kw}[1]{\texttt{#1}}
\newcommand{\bG}{\boldsymbol{\Gamma}}
\newcommand{\bE}{\mathbf{E}}
\newcommand{\bu}{\mathbf{u}}
\newcommand{\bn}{\mathbf{n}}
\newcommand{\kB}{k_\mathrm{B}}
\newcommand{\Te}{T_\mathrm{e}}
\newcommand{\nel}{n_\mathrm{e}}
\newcommand{\vth}{\bar v_\mathrm{e}}
\newcommand{\pd}[2]{\frac{\partial #1}{\partial #2}}
\newcolumntype{L}[1]{>{\raggedright\arraybackslash}p{#1}}

\title{SOMAFOAM\\[4pt]\large Model equations, numerical schemes\\ and boundary conditions}
\author{}
\date{October 2026}

\emergencystretch=3em
\begin{document}
\maketitle

\noindent
This document describes what the solver \kw{somaFoam} solves and how, as
implemented in \kw{src/plasmaCookBook} and \kw{applications/somaFoam}.

\tableofcontents

%------------------------------------------------------------------------------
\section{Overview}

SOMAFOAM is a finite-volume fluid model of low-temperature plasmas built on
foam-extend~4.0. For each time step it solves, in this order,
\begin{enumerate}
\item the Poisson equation for the potential $\Phi$ (with the dielectric
      regions, if any),
\item a continuity equation for every transported species, with a
      drift--diffusion flux or with a momentum equation,
\item the ion and gas temperatures, according to the selected models,
\item the electron energy equation for the electron temperature $\Te$.
\end{enumerate}
Temperatures are in kelvin throughout, densities in m$^{-3}$, and all other
quantities in SI units.

\begin{table}[h]
\centering
\begin{tabular}{@{}ll@{}}
\toprule
Symbol & Meaning \\
\midrule
$n_k$, $\bG_k$ & number density and particle flux of species $k$ \\
$z_k$ & charge number ($-1$ for electrons) \\
$\mu_k$, $D_k$ & mobility (positive) and diffusion coefficient \\
$T_k$ & temperature of species $k$ ($\Te$ for electrons, $T_\mathrm{i}$ for ions, $T$ for the gas) \\
$m_k$, $M$ & particle mass of species $k$ and of the background gas \\
$\nu_k$ & momentum-transfer collision frequency \\
$S_k$ & net production rate by the reactions \\
$\Phi$, $\bE=-\nabla\Phi$ & potential and electric field \\
$\varepsilon$ & permittivity \\
$\sigma$ & surface charge density on a dielectric surface \\
$\bn$ & unit normal on a boundary, pointing out of the plasma \\
$\delta$ & distance from the centre of the cell next to a boundary to the boundary face \\
\bottomrule
\end{tabular}
\end{table}

%------------------------------------------------------------------------------
\section{Governing equations}

\subsection{Species: drift--diffusion}

For species with \kw{transportModel driftDiffusion} (models \kw{driftDiffusion}
and \kw{mixed}),
\begin{equation}
\pd{n_k}{t} + \nabla\cdot\bG_k = S_k ,
\qquad
\bG_k = \mu_k\Big(z_k\bE - \frac{\kB}{e}\nabla T_k\Big)\,n_k - D_k\nabla n_k .
\label{eq:dd}
\end{equation}
The first term of the flux is the drift in the field together with the
thermal-diffusion drift; the code calls the bracket times $\mu_k$ the drift
velocity $\mathbf{F}_k$. Neutral species have $\bG_k=-D_k\nabla n_k$.

The mobility and the diffusion coefficient of each species are set by
\kw{mobilityModel} and \kw{diffusionModel} in \kw{constant/plasmaProperties}:
\begin{center}
\begin{tabular}{@{}lL{10.5cm}@{}}
\toprule
Entry & Meaning \\
\midrule
\kw{constant} & constant value given in the species dictionary \\
\kw{eTemp} & table of $\mu N$ or $DN$ against $\Te$ (\kw{constant/mu\_<species>}, \kw{D\_<species>}), divided by the gas number density $N$ \\
\kw{EON} & the same, tabulated against the reduced field $E/N$ \\
\kw{einsteinRelation} & (diffusion only) $D_k=\mu_k \kB T_k/e$ \\
\bottomrule
\end{tabular}
\end{center}

\subsection{Species: momentum equation}

For species with \kw{transportModel momentum} (models \kw{momentum} and
\kw{mixed}) the velocity $\bu_k$ is solved for:
\begin{align}
\pd{n_k}{t} + \nabla\cdot(n_k\bu_k) &= S_k , \\
\pd{(n_k\bu_k)}{t} + \nabla\cdot(n_k\bu_k\bu_k) + \nu_k n_k\bu_k
  &= \frac{1}{m_k}\Big[-\nabla(n_k\kB T_k) + z_k e\,n_k\bE\Big] .
\end{align}
The collision frequency is $\nu_k=e/(\mu_k m_k)$ with
\kw{collisionFrequency muBased}, or is taken from the reaction rates with
\kw{rrBased}.

\subsection{Reactions}

The reactions are listed in \kw{constant/reactions}. Rate coefficients are of
Arrhenius form or tabulated against $\Te$ (files \kw{constant/reaction\_*});
electron-impact reactions carry the energy lost by the electron per event.
The net rate $S_k$ of each species is linearised in its own density,
\begin{equation}
S_k \simeq S_k^{*} + \Big(\pd{S_k}{n_k}\Big)^{*}\,(n_k-n_k^{*}) ,
\end{equation}
where $^{*}$ denotes the previous inner iteration, and the derivative is
treated implicitly.

\subsection{Poisson equation}

Selected by \kw{poissonEquationSolver} in \kw{constant/electroMagnetics}.

\paragraph{\kw{explicit}}
\begin{equation}
\nabla\cdot(\varepsilon\nabla\Phi) = -e\sum_k z_k n_k .
\end{equation}

\paragraph{\kw{semiImplicit}} The space charge at the new time level is
predicted from the species equations, and the part of the prediction that
depends on the new field is moved to the left-hand side:
\begin{equation}
\nabla\cdot\Big[\Big(\varepsilon + \Delta t\,e\sum_k |z_k|\mu_k n_k\Big)\nabla\Phi\Big]
 = -e\sum_k z_k\Big[n_k^{o} + \Delta t\Big(S_k + \nabla\cdot(D_k\nabla n_k)
   + \nabla\cdot\big(\mu_k\tfrac{\kB}{e}n_k\nabla T_k\big)\Big)\Big] ,
\label{eq:semi}
\end{equation}
with $n_k^{o}$ the density at the old time level. This removes the
dielectric-relaxation limit on the time step. A third entry,
\kw{linearized}, solves the explicit equation and reconstructs the field from
the face fluxes of the potential.

\subsection{Electron energy}

Model \kw{efullImplicit} (entry \kw{energyModel \{ electron \dots; \}}):
\begin{equation}
\pd{}{t}\Big(\tfrac32 \kB \nel \Te\Big)
 + \nabla\cdot\Big(\tfrac52 \kB \Te\,\bG_\mathrm{e}\Big)
 - \nabla\cdot(\kappa_\mathrm{e}\nabla\Te)
 = -e\,\bG_\mathrm{e}\cdot\bE - Q_\mathrm{inel} - Q_\mathrm{el} .
\label{eq:Te}
\end{equation}
\begin{itemize}
\item Conductivity: $\kappa_\mathrm{e} = \tfrac52\,\kB^2 \nel \Te/(m_\mathrm{e}\nu_\mathrm{e})$,
      which with \kw{muBased} is $\tfrac52 \kB^2\nel\Te\,\mu_\mathrm{e}/e$.
\item $Q_\mathrm{inel}$: energy lost in the electron-impact reactions, from
      the reaction list.
\item $Q_\mathrm{el} = 3\kB\,\dfrac{m_\mathrm{e}}{M}\,\nu_\mathrm{e}\,\nel(\Te-T)$:
      elastic losses to the gas (with \kw{muBased}).
\item The losses are linearised in $\Te$ in the same way as the reaction
      sources.
\end{itemize}

\paragraph{Consistent fluxes.} The electron flux in \eqref{eq:Te} is the flux
of the electron continuity equation \emph{through the cell faces}, exactly as
that equation discretises it. With $\Gamma_f$ the electron flux through face
$f$ (positive out of the cell) and $\Phi_f$ the potential interpolated to the
face, the convective term uses $\tfrac52\kB\Gamma_f$ and the heating in a cell
of volume $V$ is
\begin{equation}
-\,e\,(\bG_\mathrm{e}\cdot\bE)_P
 = e\,(\bG_\mathrm{e}\cdot\nabla\Phi)_P
 = \frac{e}{V}\sum_f \Gamma_f\,(\Phi_f-\Phi_P) ,
\end{equation}
which is the discrete form of $\nabla\cdot(\bG\Phi)-\Phi\nabla\cdot\bG$. The
energy the field delivers and the energy the electrons receive then match
cell by cell. The fluxes of the species in the cells (used for the currents)
are reconstructed from the face fluxes. For electrons with the momentum model
the flux $\nel\bu_\mathrm{e}$ at the cell centres is used instead.

\subsection{Ion and gas temperature}

\begin{center}
\begin{tabular}{@{}llL{9cm}@{}}
\toprule
Entry & Model & Equation \\
\midrule
\kw{ion} & \kw{gasTemp} & $T_\mathrm{i}=T$ \\
 & \kw{gasCorrectedTemp} & $T_\mathrm{i}=T+\dfrac{M_\mathrm{i}+M}{5M_\mathrm{i}+3M}\,\dfrac{M(\mu_\mathrm{i}|\bE|)^2}{\kB}$ (Wannier) \\[6pt]
\kw{gas} & \kw{none} & $T$ constant \\
 & \kw{gfullImplicit} & steady conduction, $-\nabla\cdot(\kappa\nabla T)=\sum_\mathrm{ions} e\,\bG_k\cdot\bE + {}$ion--gas exchange \\
\bottomrule
\end{tabular}
\end{center}

\subsection{Dielectric regions}

With \kw{solutionDomain plasmaDielectric} each dielectric region solves
$\nabla\cdot(\varepsilon_\mathrm{d}\nabla\Phi)=0$. Charge accumulates on the
interface,
\begin{equation}
\pd{\sigma}{t} = \mathbf{J}\cdot\bn ,\qquad \mathbf{J}=e\sum_k z_k\bG_k ,
\label{eq:sigma}
\end{equation}
and the potential satisfies
$\varepsilon_\mathrm{d}\,\partial_n\Phi_\mathrm{d}-\varepsilon\,\partial_n\Phi=\sigma$
there (see Section~\ref{sec:bcphi}).

%------------------------------------------------------------------------------
\section{Discretisation and available schemes}

\subsection{Flux schemes for drift--diffusion}

\kw{fluxScheme} in \kw{constant/plasmaProperties} selects how the
drift--diffusion flux \eqref{eq:dd} and the convection--conduction flux of
\eqref{eq:Te} are discretised. It can also be given inside the dictionary of
a species, for that species alone.

\paragraph{\kw{fvSchemes} (default).} Drift and diffusion are separate terms,
discretised with the convection and Laplacian schemes named in
\kw{system/fvSchemes} (Table~\ref{tab:schemes}).

\paragraph{\kw{scharfetterGummel}.} Drift and diffusion are one flux. Between
the centres $P$ and $N$ of two neighbouring cells, a distance $d$ apart, with
$u$ the drift velocity normal to the face and $A$ the face area,
\begin{equation}
\Gamma_f = \frac{D A}{d}\Big[B(-\mathrm{Pe})\,n_P - B(\mathrm{Pe})\,n_N\Big],
\qquad \mathrm{Pe}=\frac{u\,d}{D},\qquad B(x)=\frac{x}{e^{x}-1} .
\end{equation}
This is the exact solution of the steady one-dimensional problem between the
two centres. It reduces to central differencing for $|\mathrm{Pe}|\ll1$ and
to upwind differencing for $|\mathrm{Pe}|\gg1$, and gives non-negative
densities without a limiter. The same form is used for the energy equation
with $u\to\tfrac52\kB\Gamma_f$ and $D\to\kappa_\mathrm{e}$. Boundary faces
other than processor and cyclic faces are treated as with \kw{fvSchemes}.
There is no non-orthogonal correction.

Written as separate terms, the scheme is central differencing for diffusion
with the convective face value $n_f=w\,n_P+(1-w)\,n_N$,
$w=1/(1-e^{-\mathrm{Pe}})-1/\mathrm{Pe}$.

\subsection{Terms and scheme entries}

\begin{longtable}{@{}L{3.5cm}L{6.3cm}L{5.4cm}@{}}
\caption{Terms, their entries in \kw{system/fvSchemes}, and the schemes that
can be used.\label{tab:schemes}}\\
\toprule
Term & Entry & Schemes \\
\midrule
\endfirsthead
\toprule
Term & Entry & Schemes \\
\midrule
\endhead
\bottomrule
\endfoot
Time derivatives & \kw{ddtSchemes default} & \kw{Euler} (first order, implicit; used in all examples) \\[3pt]
Drift of charged species, $\nabla\cdot(\mathbf{F}n)$ & \kw{div(F,Ni)} & \kw{Gauss} + \kw{upwind}, \kw{linear}, \kw{limitedLinear 1}, \kw{vanLeer}, \kw{MUSCL}, \kw{Minmod}, \kw{SuperBee}, \kw{vanAlbada}, \kw{OSPRE}, \kw{UMIST}, \kw{QUICK}, \kw{Gamma}, \kw{limitedCubic}; or \kw{fluxScheme scharfetterGummel} \\[3pt]
Diffusion of charged species & \kw{laplacian(D,Ni)} & \kw{Gauss linear} + \kw{orthogonal}, \kw{corrected}, \kw{limited} $\psi$, \kw{uncorrected}; not used with \kw{scharfetterGummel} \\[3pt]
Diffusion of neutral species & \kw{laplacian(D,Nin)} & as above \\[3pt]
Convection of electron energy & \kw{div((interpolate(eeFlux)\&S),Te)} & as \kw{div(F,Ni)}; not used with \kw{scharfetterGummel} for the electrons \\[3pt]
Electron heat conduction & \kw{laplacian(eC,Te)} & as \kw{laplacian(D,Ni)} \\[3pt]
Poisson equation & \kw{laplacian(eps,Phi)} & \kw{Gauss linear} + \kw{orthogonal} or \kw{corrected} \\[3pt]
Thermal-diffusion term in \eqref{eq:semi} & \kw{laplacian(D,T)} & as above \\[3pt]
Momentum model: convection of density and velocity & \kw{div(U,Ni)}, \kw{div(flux,Ui\_h)} & convection schemes as \kw{div(F,Ni)} \\[3pt]
Gas temperature & \kw{laplacian(alpha,T)} & Laplacian schemes \\[3pt]
Gradients ($\bE=-\nabla\Phi$, $\nabla T$) & \kw{gradSchemes default} & \kw{Gauss linear} \\
\end{longtable}

\noindent Notes.
\begin{itemize}
\item \kw{limitedLinear 1} is the scheme of the examples. With consistent
      energy fluxes it and \kw{vanLeer} give results within a few per cent of
      Scharfetter--Gummel on the meshes tested (1~Torr argon, 200 cells).
\item On meshes that are not orthogonal, including the fine/coarse interfaces
      of adaptively refined 2D and 3D meshes, use \kw{corrected} Laplacian
      and \kw{snGrad} schemes.
\item On processor boundaries the face values are interpolated between the
      two cells as on internal faces, so parallel and serial results agree.
\end{itemize}

\subsection{Time stepping and inner iterations}

One PIMPLE iteration per time step is the normal setting. Within a time step
each species equation and the energy equation are solved repeatedly to
resolve the linearised sources, as set in \kw{constant/plasmaProperties}:
\begin{center}
\begin{tabular}{@{}lll@{}}
\toprule
\kw{innerIterations} entry & Default & Meaning \\
\midrule
\kw{chargedSpecies} & 4 & passes per charged species \\
\kw{neutralSpecies} & 2 & passes per neutral species \\
\kw{electronTemperature} & 6 & passes of the energy equation \\
\kw{tolerance} & $10^{-5}$ & initial residual that ends the passes \\
\kw{reportUnconverged} & yes & report loops that end above the tolerance \\
\bottomrule
\end{tabular}
\end{center}
The time step is fixed when \kw{deltaTMin} and \kw{deltaTMax} in
\kw{system/controlDict} are equal; otherwise it follows
$\Delta t=\kw{MaxCo}/\max(\tfrac12|\nabla\cdot\mathbf{F}_\mathrm{e}|,\dots)$
between those bounds. Field relaxation factors \kw{<field>Final} in
\kw{system/fvSolution} apply on the last PIMPLE iteration.

\subsection{Limits}

Floors and ceilings are applied after each solve (\kw{limits} in
\kw{constant/plasmaProperties}; they are read again when the file changes
during a run):
\kw{densityFloor} ($10^{4}$~m$^{-3}$), \kw{TeMin} (300~K), \kw{TeMax}
(none), and per field \kw{N\_<species> \{ min; max; \}}, \kw{Te}, \kw{Tion},
\kw{T}.

\subsection{Acceleration of slow neutral species}

With \kw{slowSpeciesAcceleration}, after every \kw{fullCycles} periods of the
applied voltage the listed neutral species are advanced alone over a long
time $\Delta t_\mathrm{a}$ per step:
\begin{equation}
\frac{n_k^{\mathrm{new}}-n_k}{\Delta t_\mathrm{a}}
 - \nabla\cdot(D_k\nabla n_k^{\mathrm{new}})
 = \langle S_k\rangle - \Big\langle\pd{S_k}{n_k}\Big\rangle n_k
   + \Big\langle\pd{S_k}{n_k}\Big\rangle n_k^{\mathrm{new}} ,
\end{equation}
where $\langle\cdot\rangle$ is the average over the last \kw{averageCycles}
periods. With \kw{steadyState yes} the time derivative is dropped. The change
per advance is limited by \kw{maxChangeFactor} and \kw{relaxation}.

\subsection{Adaptive mesh refinement}

Cells are refined where an indicator exceeds \kw{lowerRefineLevel} and merged
where it is below \kw{unrefineLevel}. The indicator of a cell is the maximum
over the listed entries:
\begin{center}
\begin{tabular}{@{}ll@{}}
\toprule
\kw{type} & Value \\
\midrule
\kw{relativeGradient} & $h\,|\nabla f|\,/\,\big(|f|+\kw{floor}\big)$ \\
\kw{gradient} & $|f_1-f_2|/\kw{scale}$ between neighbouring cells 1 and 2 \\
\kw{magnitude} & $|f|/\kw{scale}$ \\
\bottomrule
\end{tabular}
\end{center}
Here $h$ is the size of the cell. With $L=|f|/|\nabla f|$ the local scale
length of the field, \kw{relativeGradient} is $h/L$: a cell is refined when
it is larger than \kw{lowerRefineLevel} times the scale length, and each
level halves $h$. The \kw{floor} keeps regions where $f$ is negligible from
being refined.

1D meshes use \kw{dynamicRefine1DFvMesh}; 2D and 3D meshes use
\kw{dynamicPolyRefinementFvMesh} with the \kw{plasmaIndicatorRefinement}
selection. Merged cells take the volume average of the fields.

%------------------------------------------------------------------------------
\section{Boundary conditions}

On a boundary face, $n_b$ and $T_b$ are the face values and $n_P$, $T_P$ the
values in the adjacent cell. Several conditions are of mixed type,
\begin{equation}
\phi_b = f\,\phi_\mathrm{ref} + (1-f)\Big(\phi_P + \delta\,g_\mathrm{ref}\Big),
\end{equation}
with weight $f$, reference value $\phi_\mathrm{ref}$ and reference gradient
$g_\mathrm{ref}$. $E_n=\bE\cdot\bn$ is positive when the field points into
the wall, that is, when electrons are repelled from it.

\subsection{Potential}\label{sec:bcphi}

\begin{longtable}{@{}L{4.9cm}L{10.3cm}@{}}
\toprule
Type & Condition \\
\midrule
\endhead
\bottomrule
\endfoot
\kw{fixedValue} & $\Phi=\text{value}$ (grounded electrode) \\[3pt]
\kw{plasmaPotential} & driven electrode, by \kw{model}:\newline
  \kw{directCurrent}: $\Phi=A$;\newline
  \kw{cosFrequencyModulated}: $\Phi=A\cos(2\pi f t)+V_\mathrm{bias}$;\newline
  \kw{sinFrequencyModulated}: $\Phi=A\sin(2\pi f t)+V_\mathrm{bias}$\newline
  (entries \kw{amplitude}, \kw{frequency}, \kw{bias}) \\[3pt]
\kw{externalCircuitPotential} & electrode behind a series resistance $R$ and capacitance $C$:
  $\Phi=A\sin(2\pi f t)+V_\mathrm{bias}+I R+Q_C/C$, with $I$ the total
  current through the patch and $Q_C=\int I\,dt$ (entries \kw{R}, \kw{C}).
  With \kw{directCurrent} the source is $A\,(1-e^{-t/\tau})$, with
  $\tau$ = \kw{rampTime} (no ramp if absent) \\[3pt]
\kw{rfPowerElectrode} & RF electrode at a given mean power behind a
  blocking capacitor: $\Phi=A\sin(2\pi f t)+V_\mathrm{dc}$. Every period the
  mean current $\langle I\rangle$ and mean power
  $\langle \Phi I\rangle$ through the patch are evaluated ($I$ is the total
  current into the plasma). Every \kw{updateCycles} periods
  $A\leftarrow A\,(P/\langle\Phi I\rangle)^{\gamma/2}$, with $P$ = \kw{power},
  $\gamma$ = \kw{relaxation}, limited to the fraction \kw{maxChange}.
  \kw{blockingCapacitor}: \kw{ideal},
  $V_\mathrm{dc}\leftarrow V_\mathrm{dc}-\langle I\rangle/(fC)$ every period
  ($C$ = \kw{capacitance} acts as a gain); \kw{physical},
  $V_\mathrm{dc}=V_\mathrm{bias}-Q/C$ with $Q=\int I\,dt$; \kw{none},
  fixed \kw{bias} \\[3pt]
\kw{dielectricSideWall} & charging wall without a meshed dielectric:
  $\partial_n\Phi=\sigma/\varepsilon_0$ (no field inside the wall), with
  $\sigma$ from \eqref{eq:sigma} \\[3pt]
\kw{coupledPotential} & plasma--dielectric interface, \kw{somaFoam}: the two
  regions are solved in one matrix; the interface potential is
  $\Phi_w=\dfrac{\varepsilon_1\Phi_1/\delta_1+\varepsilon_2\Phi_2/\delta_2+\sigma}
  {\varepsilon_1/\delta_1+\varepsilon_2/\delta_2}$ \\[3pt]
\kw{iterativeCoupledPotential} & plasma--dielectric interface,
  \kw{plasmaMultiRegionFoam}: the plasma side fixes $\Phi_w$, the dielectric
  side imposes
  $\partial_n\Phi_\mathrm{d}=(\sigma-\varepsilon\,\partial_n\Phi)/\varepsilon_\mathrm{d}$
  (normals out of each region), and $\Phi_w$ is relaxed (Aitken) until both
  agree \\
\end{longtable}

\subsection{Electron density: \kw{driftDiffusionElectronDensity}}

The condition imposes the net electron flux to the wall: the thermal flux,
always, plus the drift where it is directed to the wall, minus the emitted
electrons,
\begin{equation}
\bG_\mathrm{e}\cdot\bn = \tfrac14\vth\,n_b + \max(\mathbf{F}_\mathrm{e}\cdot\bn,\,0)\,n_b
 - \Gamma_\mathrm{se} - \Gamma_\mathrm{FE},
\qquad \vth=\sqrt{\frac{8\kB\Te}{\pi m_\mathrm{e}}} ,
\label{eq:bce}
\end{equation}
where $\mathbf{F}_\mathrm{e}\cdot\bn=-\mu_\mathrm{e}\big(E_n+\tfrac{\kB}{e}\partial_n\Te\big)$
is the drift velocity towards the wall. It is implemented as a mixed
condition with
\begin{equation}
f = \frac{C_1}{C_1+D_\mathrm{e}/\delta},\qquad
C_1=\tfrac14\vth+\max(-\mathbf{F}_\mathrm{e}\cdot\bn,\,0),
\qquad n_\mathrm{ref}=0,\qquad
g_\mathrm{ref}=\frac{\Gamma_\mathrm{se}+\Gamma_\mathrm{FE}}{D_\mathrm{e}} .
\end{equation}
$C_1$ is never negative, so $0\le f<1$ for any wall cell. With
\kw{Edepend false} the flux is the thermal flux alone, or the drift to the
wall where that is larger: $C_1=\max(\tfrac14\vth-\mathbf{F}_\mathrm{e}\cdot\bn,\,0)$.
\begin{itemize}
\item Secondary emission: $\Gamma_\mathrm{se}=\gamma\,\max(\bG_\mathrm{i}\cdot\bn,0)$,
      with $\gamma$ = \kw{seec} and $\bG_\mathrm{i}$ the total ion flux.
\item Field emission (\kw{field\_emission true}): the Fowler--Nordheim
      current density for the field $\beta E_n$ (\kw{field\_enhancement\_factor}
      $\beta$, \kw{work\_function} in eV), where $E_n>0$.
\end{itemize}

\subsection{Ion density}

\begin{center}
\begin{tabular}{@{}L{6.3cm}L{8.9cm}@{}}
\toprule
Type & Condition \\
\midrule
\kw{zeroGradient} & ions leave with their drift flux, $\bG_\mathrm{i}\cdot\bn=\mu_\mathrm{i}E_n n_b$ (used in the examples) \\[3pt]
\kw{driftDiffusionPositiveIonDensity} & as above where $E_n\ge0$; zero net flux where the field points away from the wall: mixed with $f=C_1/(C_1+D_\mathrm{i}/\delta)$, $C_1=-\mu_\mathrm{i}E_n$ \\
\bottomrule
\end{tabular}
\end{center}
Neutral species use \kw{zeroGradient} (no wall loss) or \kw{fixedValue}.

\subsection{Electron temperature: \kw{electronTemperature}}

The electron energy flux to the wall is the enthalpy carried by the
electrons (Lymberopoulos and Economou, J.~Phys.~D \textbf{28} (1995) 727),
\begin{equation}
\mathbf{q}_\mathrm{e}\cdot\bn = \tfrac52\kB\Te\,\Gamma_\mathrm{out}
 - \tfrac52\kB\big(\Gamma_\mathrm{se}T_\mathrm{se}+\Gamma_\mathrm{FE}T_\mathrm{FN}\big),
\end{equation}
with $\Gamma_\mathrm{out}$ the flux of plasma electrons to the wall. The
energy equation carries
$\tfrac52\kB\Te\,\bG_\mathrm{e}\cdot\bn-\kappa_\mathrm{e}\partial_n\Te$
through the wall face, and
$\bG_\mathrm{e}\cdot\bn=\Gamma_\mathrm{out}-\Gamma_\mathrm{se}-\Gamma_\mathrm{FE}$
(\ref{eq:bce}), so the condition on the temperature is
\begin{equation}
\kappa_\mathrm{e}\,\partial_n\Te = -h\,(T_b-T_\mathrm{em}),\qquad
h=\tfrac52\kB(\Gamma_\mathrm{se}+\Gamma_\mathrm{FE}),\qquad
T_\mathrm{em}=\frac{\Gamma_\mathrm{se}T_\mathrm{se}+\Gamma_\mathrm{FE}T_\mathrm{FN}}
{\Gamma_\mathrm{se}+\Gamma_\mathrm{FE}} .
\end{equation}
It is a mixed condition with
\begin{equation}
f=\frac{h}{h+\kappa_\mathrm{e}/\delta},\qquad T_\mathrm{ref}=T_\mathrm{em},\qquad g_\mathrm{ref}=0 .
\end{equation}
$h$ is never negative, so $0\le f<1$ for any size of the cell next to the
wall; without emission the condition is a zero gradient. Entries:
\kw{seec}, \kw{Tse}, \kw{TFN}, \kw{field\_emission},
\kw{field\_enhancement\_factor}, \kw{work\_function} (\kw{Edepend} is
accepted and not used).

The value $\tfrac52\kB\Te$ per electron is the one consistent with the
convective term of \eqref{eq:Te}. With the one-way kinetic value
$2\kB\Te$ the difference would have to be returned by conduction with a
gradient proportional to the wall temperature, which is unstable on coarse
wall cells.

\subsection{Auxiliary fields}

The files \kw{0/eeFlux} and \kw{0/F\_electron} hold boundary values of the
thermal energy flux and thermal velocity (\kw{eEnergyThermalFlux},
\kw{electronThermalFlux}); they are needed for the fields to be created, and
are used by the energy equation only for electrons with the momentum model.

%------------------------------------------------------------------------------
\section{Diagnostics}

The \kw{electrodeVoltageCurrent} function object writes, per patch,
\begin{equation}
I_\mathrm{cond}=\int \mathbf{J}\cdot d\mathbf{S},\qquad
I_\mathrm{disp}=\int \varepsilon\,\pd{\bE}{t}\cdot d\mathbf{S},\qquad
I_\mathrm{tot}=I_\mathrm{cond}+I_\mathrm{disp},
\end{equation}
with the sign reversed so that a current is positive from the electrode into
the plasma. The \kw{fieldAverage} function object with
\kw{resetOnOutput true} and a write interval of whole periods gives
period-averaged fields.

\end{document}
